\documentclass[11pt, a4paper]{article}
\usepackage[utf8]{inputenc}
\usepackage[T1]{fontenc}
\usepackage{amsmath, amssymb, amsthm}
\usepackage{geometry}
\usepackage{enumitem}
\usepackage{xcolor}
\usepackage{tcolorbox}

\geometry{margin=2cm}
\setlength{\parindent}{0pt}
\setlength{\parskip}{0.6em}

\definecolor{tumblue}{HTML}{0065BD}
\definecolor{tumdark}{HTML}{003359}
\definecolor{accent}{HTML}{E37222}

\newcommand{\actiontitle}[1]{{\Large\bfseries\color{tumdark} #1}\par\vspace{0.3em}}
\newcommand{\oldtitle}[1]{{\small\color{gray} Alt: #1}\par}
\newcommand{\slidenote}[1]{{\small\color{accent}\textbf{Hinweis:} #1}\par}

\title{Revidierte Folien --- Bottom-Up Herleitung\\[0.3em]
\large Akt~1: Von Entscheidungssystemen zu Boundary Testing}
\author{Arbeitsversion f\"ur Ruben}
\date{April 2026}

\begin{document}
\maketitle

Jede Folie unten enth\"alt: neuen Action Title, Body-Text, LaTeX-Formeln zum lokalen Rendern, und Kommentare zur \"Anderung gegen\"uber der alten Version.

\bigskip
\hrule
\bigskip

%% ===== SLIDE 1 (ersetzt alte A1: "Every AI testing paper...") =====
\section*{Folie 1 --- Einstieg}

\actiontitle{A decision system maps inputs to categorical outputs.}

\oldtitle{Every AI testing paper starts with the same argument.}

\textbf{Body:}

Any system that receives an input and returns a categorical output is a decision system:
%
\begin{tcolorbox}[colframe=tumblue, colback=white]
\begin{equation}
f \colon M \to Y, \qquad |Y| \geq 2
\end{equation}
\end{tcolorbox}

Three examples, same structure:

\begin{center}
\begin{tabular}{lll}
\textbf{System} & \textbf{Input} $\boldsymbol{m \in M}$ & \textbf{Output} $\boldsymbol{y \in Y}$ \\
\hline
Radiologist & Chest X-ray & Diagnosis \\
VLM (Qwen) & Image + Prompt & Classification \\
Credit Scoring & Application & Approve / Reject \\
\end{tabular}
\end{center}

\textit{All three share the same formal structure. What follows is general.}

\slidenote{Kein Zitat mehr als Einstieg. Wir starten direkt mit der Struktur. Das DeepXplore-Zitat kann als Footnote auf einer sp\"ateren Folie erscheinen, um den Bezug zur Literatur herzustellen.}

\bigskip\hrule\bigskip

%% ===== SLIDE 2 (ersetzt alte A2+A3, zusammengefasst) =====
\section*{Folie 2 --- Input-Raum}

\actiontitle{The input space is structured and context-dependent.}

\oldtitle{Every system that decides creates boundaries. / The context determines the decision space.}

\textbf{Body:}

Each input dimension $j$ has a value space $S_j$:
%
\begin{equation}
S_j \in \bigl\{\mathbb{R},\; \mathbb{Z},\; \{c_1, \ldots, c_k\}\bigr\}
\end{equation}

The combinatorially complete space is:
%
\begin{equation}
U = \prod_{j=1}^{d} S_j
\end{equation}

In a given context $C$ (dataset, prompt, deployment setting), only a subset is realized:
%
\begin{tcolorbox}[colframe=tumblue, colback=white]
\begin{equation}
M = \sigma(C) \subseteq U
\end{equation}
\end{tcolorbox}

\textit{Example:} A 2-class prompt and a 15-class prompt over the same images define different $C$, therefore different $M$, therefore different boundary structures.

\slidenote{Zwei alte Folien zu einer verschmolzen. Prompt-Konfiguration als konkretes Beispiel statt abstrakter Behauptung. Bottom-up: erst der Raum, dann die Struktur darin.}

\bigskip\hrule\bigskip

%% ===== SLIDE 3 (ersetzt alte A2, zweite H\"alfte) =====
\section*{Folie 3 --- Partition und Boundaries}

\actiontitle{Choices partition $M$ into regions --- boundaries follow.}

\textbf{Body:}

$f$ assigns each $m \in M$ to exactly one $y_k \in Y$. This creates a partition:
%
\begin{equation}
R_k = \bigl\{m \in M \;\big|\; f(m) = y_k\bigr\}, \qquad M = \bigsqcup_{k=1}^{K} R_k
\end{equation}

Where two regions meet, a \textbf{boundary} exists. This is not an assumption --- it is a structural consequence of $|Y| \geq 2$:
%
\begin{equation}
\partial R_{j} \cap \partial R_{k} \neq \emptyset \quad \text{for at least one pair } (j, k)
\end{equation}

\textit{The question is not \emph{whether} boundaries exist, but what their geometry looks like.}

\slidenote{Eigene Folie f\"ur die Partition. Die Boundary entsteht hier als logische Konsequenz, nicht als Behauptung. Kein Psychologie-Argument n\"otig.}

\bigskip\hrule\bigskip

%% ===== SLIDE 4 (ersetzt alte B1: contrast function) =====
\section*{Folie 4 --- Die Kontrastfunktion $g$}

\actiontitle{The contrast function $g$ makes boundaries measurable.}

\oldtitle{The contrast function g assigns each input a scalar.}

\textbf{Body:}

For a class pair $(j, k)$, define:
%
\begin{tcolorbox}[colframe=tumblue, colback=white]
\begin{equation}
g_{jk}(m) = P(y_j \mid m) - P(y_k \mid m)
\end{equation}
\end{tcolorbox}

\begin{align}
g_{jk}(m) > 0  &\quad \Rightarrow \quad \text{system prefers } y_j \\
g_{jk}(m) < 0  &\quad \Rightarrow \quad \text{system prefers } y_k \\
g_{jk}(m) = 0  &\quad \Rightarrow \quad \text{boundary}
\end{align}

Properties:
\begin{itemize}[nosep]
\item $|g_{jk}(m)|$ is a continuous measure of boundary proximity
\item Antisymmetry: $g_{jk} = -g_{kj}$
\item Boundary as derived object: $B_{jk} = g^{-1}(0) = \{m \in M \mid g_{jk}(m) = 0\}$
\end{itemize}

$g$ is the primary formal object. $B$ is derivable from $g$, but $g$ carries more information.

\slidenote{Inhaltlich gleich, aber die Folie kommt jetzt nach der Partition-Folie, also bottom-up: erst Regionen, dann das Werkzeug, das sie vermisst.}

\bigskip\hrule\bigskip

%% ===== SLIDE 5 (ersetzt alte B2: four quantities) =====
\section*{Folie 5 --- Vier geometrische Gr\"o\ss en}

\actiontitle{Four geometric quantities follow from $g$ and $\Delta g$.}

\textbf{Body:}

On the genotype lattice $G \subset \mathbb{Z}^d$, define the lattice difference:
%
\begin{equation}
\Delta_i\, g(x) = g(x + e_i) - g(x)
\end{equation}

This is exact --- not an approximation. Four quantities emerge:

\begin{center}
\renewcommand{\arraystretch}{1.4}
\begin{tabular}{llp{5.5cm}}
\textbf{Quantity} & \textbf{Definition} & \textbf{Measures} \\
\hline
Margin & $|g_{jk}(m)|$ & Distance to boundary in output space \\
Sensitivity & $\|\Delta g(m)\| = \max_i |\Delta_i g|$ & How much $g$ changes per unit perturbation \\
Thickness & $\dfrac{|g|}{\|\Delta g\|}$ & Perturbation budget until boundary reached \\
Direction & $\arg\max_i |\Delta_i g(m)|$ & Most decision-relevant input dimension \\
\end{tabular}
\end{center}

\slidenote{Gleiche Tabelle, aber jetzt mit der Lattice-Definition vorangestellt statt als Footnote. Macht den diskreten Ansatz explizit.}

\bigskip\hrule\bigskip

%% ===== SLIDE 6 (ersetzt alte A4+A5+A6+A7: drei Facetten → eine Folie) =====
\section*{Folie 6 --- Praktische Relevanz (eine Folie statt vier)}

\actiontitle{Each quantity answers a different question about the system.}

\oldtitle{Without boundary knowledge, the system is unpredictable / uncontrollable / unevaluable.}

\textbf{Body:}

\begin{center}
\renewcommand{\arraystretch}{1.6}
\begin{tabular}{lll}
\textbf{Question} & \textbf{Quantity} & \textbf{Formal} \\
\hline
``How clear was this decision?'' & Margin & $|g_{jk}(m)|$ \\
``How fragile is this prediction?'' & Thickness & $|g| \,/\, \|\Delta g\|$ \\
``Which input dimension matters?'' & Direction & $\arg\max_i |\Delta_i g|$ \\
``How responsive is $g$ here?'' & Sensitivity & $\|\Delta g(m)\|$ \\
\end{tabular}
\end{center}

All four derived from the same $g_{jk}$.

\begin{tcolorbox}[colframe=gray, colback=white, fontupper=\small]
\textbf{Side note:} These questions map independently to established concepts in systems psychology --- competence need (Deci \& Ryan, 1985), perceived control (Bandura, 1989), and procedural justice (Thibaut \& Walker, 1975). The exact scope of this interdisciplinary connection is part of ongoing research.
\end{tcolorbox}

\slidenote{Vier alte Folien (Predictability, Controllability, Evaluability, Summary) zu einer zusammengefasst. Psychologie ist jetzt eine Footnote/Side-Box, nicht das tragende Argument. Die Herleitung steht auf eigenen formalen Beinen.}

\bigskip\hrule\bigskip

%% ===== SLIDE 7 (neue Folie: Bezug zur Literatur) =====
\section*{Folie 7 --- Literatur-Alignment}

\actiontitle{Existing work operates on the same structure --- without naming it.}

\textbf{Body:}

\begin{quote}
\textit{``Deep learning systems are increasingly deployed in safety- and security-critical domains [\ldots] where the correctness and predictability of a system's behavior for corner case inputs are of great importance.''}

\hfill --- Pei et al.\ (DeepXplore), SOSP 2017
\end{quote}

In our notation:
\begin{itemize}[nosep]
\item ``corner case inputs'' $\longleftrightarrow$ inputs near $B_{jk}$ (low margin)
\item ``correctness'' $\longleftrightarrow$ whether $f(m)$ is stable (thickness)
\item ``predictability'' $\longleftrightarrow$ knowing $g(m)$ before observing $f(m)$
\end{itemize}

The framework gives these intuitions a shared formal object: $g_{jk}$.

\slidenote{Das DeepXplore-Zitat kommt jetzt \emph{nach} der Formalisierung als Br\"ucke, nicht als Er\"offnung. Wir zeigen: die Literatur redet \"uber dasselbe, nur ohne die gemeinsame Sprache. Kein Vorwurf (``uncreative''), sondern Einordnung.}

\bigskip\hrule\bigskip

%% ===== SLIDE 8 (ersetzt alte B4: Opacity → BT) =====
\section*{Folie 8 --- Boundary Testing}

\actiontitle{For opaque systems, $g$ is only empirically accessible.}

\textbf{Body:}

\begin{align}
\textbf{Transparent } f &\quad \Rightarrow \quad g \text{ approximable through inspection} \\
\textbf{Opaque } f \text{ (DNN, VLM)} &\quad \Rightarrow \quad g \text{ only accessible through queries}
\end{align}

\begin{tcolorbox}[colframe=accent, colback=white]
\begin{equation}
BT(f) \;:=\; \text{systematic empirical approximation of } g_{jk}
\end{equation}
\end{tcolorbox}

Complementary to XAI:
\begin{center}
\begin{tabular}{ll}
XAI (LIME, SHAP) & $\to$ local, post-hoc, approximates $g$ at single points \\
BT & $\to$ global, systematic, maps the $g$-landscape \\
\end{tabular}
\end{center}

VLMs are the hardest case --- but not the reason for the framework.

\slidenote{Inhaltlich fast gleich. Positionierung ist jetzt der logische Schluss der Bottom-Up-Kette: Raum $\to$ Partition $\to$ $g$ $\to$ Quantities $\to$ Opacity $\to$ BT.}

\bigskip\hrule\bigskip

%% ===== \"UBERSICHT =====
\section*{Zusammenfassung: Neue Folienreihenfolge}

\begin{center}
\renewcommand{\arraystretch}{1.3}
\begin{tabular}{clp{6cm}}
\textbf{\#} & \textbf{Action Title} & \textbf{Kern} \\
\hline
1 & A decision system maps inputs to categorical outputs. & $f: M \to Y$, drei Beispiele \\
2 & The input space is structured and context-dependent. & $S_j$, $U$, $M = \sigma(C)$, Prompt als $C$ \\
3 & Choices partition $M$ into regions --- boundaries follow. & $R_k$, Partition, $\partial R_j \cap \partial R_k$ \\
4 & The contrast function $g$ makes boundaries measurable. & $g_{jk}(m)$, Vorzeichen, $B_{jk}$ \\
5 & Four geometric quantities follow from $g$ and $\Delta g$. & Tabelle, Lattice-Definition \\
6 & Each quantity answers a different question about the system. & Fragen $\leftrightarrow$ Quantities, Psych.\ als Side-Note \\
7 & Existing work operates on the same structure --- without naming it. & DeepXplore-Zitat $\leftrightarrow$ Framework-Notation \\
8 & For opaque systems, $g$ is only empirically accessible. & BT-Definition, XAI-Abgrenzung \\
\end{tabular}
\end{center}

\textbf{Alt:} 9~Motivationsfolien + 4~Formalisierungsfolien = 13 Folien, top-down, Psychologie als Hauptargument.

\textbf{Neu:} 8~Folien, bottom-up, Formalisierung \emph{ist} die Motivation, Psychologie als einordnende Footnote.

\end{document}
